% Avsnitt 4.1-4.9, ur lectures/L04/appendix/a_metastability_and_synchronization.md (A.1-A.9).

\section{Problemet: en knapp är inte en testbänkssignal}
\label{c4:sec:problem}

I \chapref{c3:ch} byggde du flankdetektorer och resonerade om dem mot en ren klocka, där varje
ingång ändrades prydligt mellan flankerna, precis så som en simulator ritar det. En tryckknapp, en
extern resetledning, eller vilken signal som helst från en annan klockdomän ändras när omvärlden
bestämmer sig för det, helt oberoende av din FPGA:s klocka.

Varje D-vippa har två tidskrav kring den aktiva klockflanken:
\begin{itemize}
  \item en \textbf{setup-tid}: hur länge \code{D} måste vara stabil \emph{före} flanken.
  \item en \textbf{hold-tid}: hur länge \code{D} måste förbli stabil \emph{efter} den.
\end{itemize}

Om en ingång ändras inuti det fönstret är vippan inte garanterad att fånga en ren \code{0} eller
\code{1}, och du kan aldrig lova att den hamnar utanför, eftersom du inte styr när en extern
signal ändras. Resten av det här kapitlet löser det inte genom att undvika problemet utan genom
att innesluta det.

% ------------------------------------------------------------------------------------------
\section{Metastabilitet: vad som faktiskt händer inuti vippan}
\label{c4:sec:meta}

En D-vippa är byggd av korskopplade grindar som slår om till \code{0} eller \code{1} i samma
ögonblick som klockan samplar ingången. Om \code{D} ändras för nära flanken kan den interna
återkopplingsslingan lämnas i balans mellan de två tillstånden, varken helt hög eller helt låg,
längre än väntat. Det är \textbf{metastabilitet}.

Tänk dig en kula som balanserar exakt på toppen av en kulle. Med tillräckligt med tid rullar den
ner åt ena hållet, men exakt när den bestämmer sig, och åt vilket håll, går inte att förutsäga från
det ögonblick den placerades där, och så länge den fortfarande balanserar är dess position inte
något giltigt digitalt värde alls.

En metastabil utgång beter sig likadant:
\begin{itemize}
  \item den stabiliserar sig nästan alltid till en giltig \code{0} eller \code{1} av sig själv,
    givet tillräckligt med tid.
  \item ”tillräckligt med tid” är inte gratis, och är inte densamma varje klockcykel.
  \item om den ostabiliserade spänningen når logik nedströms innan den lagt sig kan olika grindar
    med all rätt läsa den olika: en ser \code{0}, en annan \code{1}, i samma ögonblick.
\end{itemize}

Den oenigheten är där riktiga system beter sig fel: en räknare som räknar upp i en del av kretsen
men inte i en annan, eller en tillståndsmaskin i ett omöjligt tillstånd. Allt i det här kapitlet
finns till för att se till att metastabiliteten stabiliserar sig \emph{innan} en signal rör resten
av din design.

I tid ser hela problemet ut så här:

\lecfig[0.72]{lectures/L04/appendix/images/setup_hold.png}{\code{D} ändras inuti
setup/hold-fönstret kring en stigande klockflank, vilket lämnar \code{Q} metastabil innan den sent
lägger sig på endera nivån.}{c4:fig:setuphold}

Läs det från vänster till höger. Det skuggade bandet är fönstret från \secref{c4:sec:problem},
$t_{su}$ före flanken och $t_h$ efter den. \code{D} ändras inuti det, vilket är precis vad en
ingång du inte styr står fritt att göra. \code{Q} är sedan varken \code{0} eller \code{1} under en
stund, och när den väl lägger sig är både \emph{när} den lägger sig och \emph{vilken nivå} den
lägger sig på utanför din kontroll. Båda de streckade fortsättningarna är legitima utfall av samma
ingång.

Lägg märke till vad diagrammet inte visar: ett felaktigt svar. Att sampla en signal mitt i en
övergång får ge dig endera nivån, och det är i sig harmlöst, eftersom ingången faktiskt höll på att
ändras i det ögonblicket. Problemet är fördröjningen innan du får ett svar över huvud taget, och
det är den nästa avsnitt köper tid för.

% ------------------------------------------------------------------------------------------
\section{Dubbelvippsynkroniseraren}
\label{c4:sec:sync}

Standardlösningen är \textbf{dubbelvippsynkroniseraren}:
\begin{itemize}
  \item Placera två D-vippor i serie på varje asynkron ingång: varje ingång som ändras utan någon
    fast relation till din klocka. En tryckknapp, en resetledning, en signal som anländer från en
    annan klockdomän.
  \item En ingång som redan klockas av \emph{denna} klocka är inte asynkron, och får \textbf{inte}
    synkroniseras. Att göra det fördröjer den bara, och förskjuter dess tidsläge relativt allt runt
    omkring. \code{seq_detect_mealy} i \chapref{c8:ch} är precis det fallet: dess \code{din} är
    synkron seriell data, så den går rakt in i maskinen medan \code{reset_n} fortfarande får en
    synkroniserare.
  \item Använd aldrig något annat än den \emph{andra} vippans utgång på andra ställen i designen.
\end{itemize}

\begin{textdrawing}
async_in
   │
   ▼
 [FF1] ──────► [FF2] ──────► synkroniserad, säker att använda
   (kan bli       (extremt sannolikt
   metastabil)     stabil)
\end{textdrawing}

Samma två vippor, i tid:

\lecfig[0.72]{lectures/L04/appendix/images/synchronizer_waveform.png}{\code{async_in} ändras strax
före en klockflank, \code{s1} blir metastabil under en klockperiod, och \code{s2} kommer ut ren en
flank senare.}{c4:fig:syncwave}

\code{FF1} är den som är exponerad för den asynkrona ingången, så den är den som kan bli
metastabil. Dess utgång behöver bara nå \code{FF2}, och den har en hel klockperiod på sig.
Sannolikheten att en metastabil nod fortfarande är ostabiliserad efter en hel period är liten, och
krymper ju längre tid den får, så när \code{FF2} samplar på nästa flank har \code{FF1} nästan
alltid lagt sig. \code{FF2}:s utgång är det du bygger resten av din synkrona logik på.

Det här är sannolikhetsbaserat, inte absolut, och \secref{c4:sec:mtbf} säger vad det innebär. Vid
de hastigheter som används i den här boken, tiotals MHz, är två steg det gängse,
industriaccepterade svaret.

En regel knyter ihop det: \textbf{varje asynkron ingång får sin egen dubbelvippkedja innan den
används någon annanstans.} Externa resetknappar, tryckknappar, och senare varje
\textbf{enbitssignal} som korsar från en annan klockdomän.

\begin{warning}
Det sista ordet är bärande, och det är den enda begränsning som är värd att lära sig innan du
överanvänder regeln. En dubbelvippkedja synkroniserar \textbf{en bit}. Kör en flerbitsbuss genom
en kedja per bit och varje bit är säker för sig men de är inte säkra \emph{tillsammans}: var och en
stabiliserar sig oberoende till det gamla eller det nya värdet på samma flank, så mottagaren kan
sampla en kombination som aldrig fanns på bussen. En räknare som stegar från \code{0111} till
\code{1000} har alla bitar ändrade på en gång, och en synkroniserare per bit kan ge dig
\code{1111} eller \code{0000}, värden räknaren aldrig höll.

Bussar som korsar klockdomäner behöver en handskakning, en FIFO, eller Gray-kodning, som ordnar det
så att bara en bit någonsin ändras åt gången. Ingenting i del~\ref{part:vhdl} korsar klockdomäner
med en buss, så du kommer inte att behöva dem där; grupprojektet återkommer till begränsningen i
\chapref{c12:ch}, där samma synkroniserare skrivs en gång och återanvänds. Att veta att
tvåvippsvaret stannar vid en bit räcker för nu.
\end{warning}

% ------------------------------------------------------------------------------------------
\section{Intuition för tidsförhållandena: varför ”nästan säkert” duger}
\label{c4:sec:mtbf}

En synkroniserare köper inte ”ingen metastabilitet”. Den köper metastabilitet som får en hel
klockperiod på sig att stabilisera sig innan den kan påverka något annat, och det är ett
statistiskt argument snarare än ett bevis. Den korrekta hållningen är ”rimligt säker, med stöd av
en extremt låg, väl karakteriserad felfrekvens”, inte ”bevisbart säker under alla förhållanden”,
vilket är skälet till att extra steg finns för extremfall i stället för att två förklaras
tillräckliga av en naturlag.

Sambandet har en standardform, värd att se en gång även om du inte kommer att räkna med den här:

\[ \mathit{MTBF} = \frac{e^{\,t_r / \tau}}{T_w \cdot f_{\mathrm{clock}} \cdot f_{\mathrm{data}}} \]

Läs den för vad den säger snarare än för siffror. $t_r$ är den stabiliseringstid ett steg får,
ungefär en klockperiod; $\tau$ och $T_w$ är båda egenskaper hos vippan, där $\tau$ är hur snabbt
ett metastabilt värde avklingar och $T_w$ bredden på det fönster kring klockflanken där ingången
över huvud taget kan framkalla ett sådant. Att lägga till ett steg köper ytterligare ett $t_r$ i
exponenten, så \textbf{täljaren växer exponentiellt med varje steg du lägger till}, mot en
klockcykels ökad latens. Nämnaren växer linjärt med klockfrekvensen och, precis lika viktigt, med
$f_{\mathrm{data}}$: hur ofta den asynkrona ingången faktiskt ändras. En knapp som trycks en gång i
sekunden och en dataledning som växlar i 10~MHz skiljer sig sju tiopotenser åt enbart i den termen,
vilket är skälet till att argumentet för ett tredje steg är en fråga om ingången, inte om klockan.

% ------------------------------------------------------------------------------------------
\section{Att synkronisera en resetsignal: aktivera asynkront, släpp synkront}
\label{c4:sec:resetsync}

Reset förtjänar ett eget mönster, som \code{reset_sync} demonstrerar. Det finns ingen
\file{reset_sync.vhd} att öppna i kursrepot: den är Övning~\ref{c4:ex:resetsync}, och du skriver
den. Listningen nedan är hela modulens process.

\begin{vhdlcode}
process(clock, reset_n) is
begin
    if (reset_n = '0') then
        reset_s1_n <= '0';
        reset_s2_n <= '0';
    elsif (rising_edge(clock)) then
        reset_s1_n <= '1';
        reset_s2_n <= reset_s1_n;
    end if;
end process;
\end{vhdlcode}

\begin{itemize}
  \item \textbf{Aktivera asynkront.} När \code{reset_n} går låg tvingas båda stegen omedelbart till
    \code{0}, utan något beroende av klockan. En reset som väntade på en klockflank skulle
    omintetgöra sitt eget syfte: den måste få verkan i samma ögonblick som den behövs, med eller
    utan klocka.
  \item \textbf{Släpp synkront.} När \code{reset_n} återgår till \code{1} försvinner inte
    resetvillkoret överallt på en gång. Det skiftas ut genom de två stegen som vilken annan
    asynkron ingång som helst, så \code{reset_s2_n}, den signal designen faktiskt använder, släpper
    rent på en klockflank tillsammans med allt annat. Utan detta skulle olika vippor kunna lämna
    reset på olika cykler och för ett ögonblick vara oense om systemets tillstånd, samma sorts fel
    som metastabilitet orsakar.
\end{itemize}

Det är därför \code{reset_n} synkroniseras av sin egen kedja innan något annat beror på den,
inklusive knappsynkroniseraren, vars resetingång är \code{reset_s2_n} snarare än den råa
\code{reset_n}.

% ------------------------------------------------------------------------------------------
\section{Att återanvända kedjan för flankdetektering, och studsfiltrering i förbigående}
\label{c4:sec:buttonsync}

Knappen behöver synkroniseras, och från \chapref{c3:ch} behöver du dessutom detektera \emph{när}
den trycks ned snarare än dess nivå. \code{button_sync}, som är Övning~\ref{c4:ex:buttonsync} och
därmed också din att skriva snarare än en fil i kursrepot, får båda ur en enda kedja av \textbf{tre}
vippor:

\begin{textdrawing}
button_n
   │
   ▼
[FF1] → [FF2] → [FF3]
        │        │
        │        └─ föregående värde
        └─ nuvarande värde
\end{textdrawing}

Vippa 1 och 2 är \secref{c4:sec:sync}:s synkroniserare, vilket gör \code{button_s2_n} till det
stabila nuvarande värdet. Vippa 3 lagrar förra cykelns \code{button_s2_n} som \code{button_s3_n},
det föregående värdet. Vippa 1 och 2 är värda att lägga märke till för sig:
Övning~\ref{c5:ex:sync} bryter ut den dubbelvippkedjan till en återanvändbar \code{sync}-modul, och
lämnar den tredje vippan här som flankdetektorns egen. Att jämföra steg 2 mot steg 3 ger
flankdetektering gratis:

\begin{vhdlcode}
button_edge_s2 <= (not button_s2_n) and button_s3_n;
\end{vhdlcode}

vilket är sant under exakt en klockcykel: den där knappen läses som nedtryckt nu och inte var det
cykeln före.

\textbf{Synkronisering löser i sig bara metastabilitet.} Den säger ingenting om en mekanisk knapps
benägenhet att studsa: att sluta och bryta kontakt flera gånger under den första millisekunden
eller så, där varje studs ser ut som ett eget tryck. Det som håller den här bokens övningar från
att dränkas i studs är \emph{kombinationen} av kedjan och ett detekteringsfönster som är långsamt i
förhållande till studsen:
\begin{itemize}
  \item I det här kapitlets CircuitVerse-övningar är klockperioden en hel sekund. Studsen lägger
    sig på långt mindre än så, så vid nästa flank läser \code{button_s2_n} en enda ren nivå. Kedjan
    ser ut att studsfiltrera, men den samplar egentligen bara för sällan för att se studsen.
  \item På FPGA:n är klockan \code{50 MHz}, ett sampel var \code{20 ns}, medan studsen varar från
    under en millisekund till några tiotals millisekunder, många tusen cykler. Vid den hastigheten
    detekterar \code{button_sync} troget \emph{varje} studsövergång som en egen flank, och
    lysdioden kan växla flera gånger för ett enda fysiskt tryck.
\end{itemize}

Så exakt: dubbelvippsynkroniseraren löser metastabilitet tillförlitligt vid de hastigheter som
används här. Om den dessutom \emph{ser ut} att studsfiltrera beror helt på att klockperioden är
långsam i förhållande till knappens studstid, vilket stämmer i CircuitVerse och inte automatiskt
stämmer på FPGA:n.

% ------------------------------------------------------------------------------------------
\section{Vad en riktig studsfiltrering i hårdvara tillför}
\label{c4:sec:debounce}

Den här boken har inget eget kapitel om studsfiltrering, och de genomarbetade exemplen fortsätter
medvetet att använda \secref{c4:sec:buttonsync}:s krets, studs och allt. Den duger för en
demolysdiod. En produktionsdesign lägger till något av följande \emph{ovanpå}:
\begin{itemize}
  \item \textbf{Ett analogt RC-lågpassfilter}, eventuellt med en Schmitt-triggerbuffert, på
    knappingången innan den når en FPGA-pinne. RC-paret jämnar ut de snabba slut- och
    brytövergångarna till en enda gradvis flank, och Schmitt-triggern, med sina två omslagströsklar,
    gör tillbaka det till en enda ren digital övergång. Det löser studsen innan signalen ens är
    digital, men kostar komponenter på kortet och hjälper inte för ingångar som inte är mekaniska
    kontakter.
  \item \textbf{En digital studsfiltreringsräknare.} I stället för att agera på nästa detekterade
    flank, kräv att ingången håller en ny nivå under någon minsta tid, några millisekunders värde
    av klockcykler, innan den accepteras. Det bygger ditt eget konstgjort långsamma
    detekteringsfönster på chippet, samma effekt som CircuitVerses 1-sekundersklocka ger gratis,
    utan att sakta ner systemklockan. Du möter räknare i \chapref{c6:ch} och timers, som den här
    tekniken behöver, i \chapref{c7:ch}.
\end{itemize}

Båda sitter ovanpå synkroniseraren, inte i stället för den: även en perfekt studsfiltrerad signal
är fortfarande asynkron mot din klocka och behöver fortfarande en dubbelvippkedja.

% ------------------------------------------------------------------------------------------
\section{Generics: en modul, flera storlekar}
\label{c4:sec:generics}

\code{reset_sync} och \code{button_sync} är byggda av samma idé, men bara den ena har en fråga att
besvara. En reset är en bit och kommer alltid att vara det. Hur många \emph{knappar} har en design?
Det här kapitlets \code{led_toggle_sync} har en, \chapref{c8:ch}:s \code{fsm_led} har två. Att
skriva ut samma trevippskedja en gång per bredd är precis den dubblering som uppdelningen i moduler
var till för att undvika.

En \textbf{generic} besvarar den: ett värde som en entitet deklarerar vid sidan av sina portar,
fastlagt när designen byggs snarare än föränderligt medan den kör.

\begin{vhdlcode}
entity button_sync is
    generic(COUNT: natural range 1 to 3 := 1);
    port(clock, reset_s2_n: in  std_logic;
         button_n         : in  std_logic_vector(COUNT-1 downto 0);
         button_edge_s2   : out std_logic_vector(COUNT-1 downto 0));
end entity;
\end{vhdlcode}

\begin{itemize}
  \item \code{COUNT} beter sig som en konstant inuti arkitekturen. Den kan inte ändras medan
    kretsen kör; den avgörs innan en enda grind placerats.
  \item \code{:= 1} är dess standardvärde, som används av varje instansiering som inte säger något.
  \item Portbredderna är skrivna \emph{i termer av den}, och så är varje intern signal som följer.
\end{itemize}

\subsection*{Att åsidosätta den}
En instansiering skickar generics med en \code{generic map}, före sin \code{port map}:

\begin{vhdlcode}
button_sync1: entity work.button_sync
    generic map(2)
    port map(clock, reset_s2_n, button_n, button_edge_s2);
\end{vhdlcode}

Generics matchas \textbf{positionellt}, i deklarationsordning, precis som portar gör i
\secref{c2:sec:submodules}. Att utelämna \code{generic map} tillämpar standardvärdet; den här boken
skriver ut den explicit även när den stämmer, eftersom en läsare inte ska behöva öppna modulen för
att ta reda på hur bred en instans är.

\subsection*{När en generic gör rätt för sig}
\code{button_sync} har en och \code{reset_sync} har ingen, och den skillnaden är hela regeln:
\textbf{lägg till en generic när något genuint varierar mellan instanser, inte av princip.}

\code{button_sync} instansieras med \code{COUNT = 1} av \secref{c4:sec:worked}:s
\code{led_toggle_sync} och med \code{COUNT = 2} av \chapref{c8:ch}:s \code{fsm_led}: två bredder,
en fil. \code{reset_sync} är en bit överallt, så en \code{WIDTH}-generic på den vore en parameter
som ingen någonsin skickar något annat än \code{1} till, och en sak till för en läsare att
kontrollera. Den andra instansieringen är det som gör en generic värd att ha; en parameter med
exakt ett värde i hela boken vore bara ett extra mellanled. \code{button_sync}s egen testbänk
instansierar den med \emph{båda} bredderna i samma simulering, just för att bevisa att en fil
verkligen betjänar båda.

Du möter ytterligare en generic i \chapref{c7:ch}: \code{timer}s \code{TICK_COUNT}, som är ett
antal snarare än en bredd, och som är det som låter samma fil mäta en femtiondels sekund i en
design och en hel sekund i en annan.

% ------------------------------------------------------------------------------------------
\section{Genomarbetat exempel: \code{led_toggle_sync}}
\label{c4:sec:worked}

Det här knyter ihop \secref{c4:sec:problem}--\ref{c4:sec:buttonsync} till \chapref{c3:ch}:s
växlingskrets, gjord säker för en riktig, asynkron knapp och reset.

\begin{booktable}{@{}p{5.6em}p{3.4em}L@{}}
  \thead{Port} & \thead{Riktn.} & \thead{Beskrivning} \\
  \midrule
  \code{clock} & in & \code{50 MHz} systemklocka på DE0-CV-kortet. \\
  \code{reset_n} & in & Aktiv låg asynkron reset, kopplad till en tryckknapp. \\
  \code{button_n} & in & Aktiv låg tryckknapp; växlar lysdioden på sin fallande flank. \\
  \code{led} & out & Lysdiod, växlas vid varje synkroniserat, flankdetekterat knapptryck. \\
\end{booktable}

Interna signaler: \code{led_s} håller lysdiodens tillstånd, \code{reset_s2_n} är reseten efter
tvåvippsynkroniseraren (suffixet \code{s2} är bokens konvention för ”synkroniserad med två
vippor”), och \code{button_edge_s2} är den encykelspuls som markerar en synkroniserad fallande
flank på \code{button_n}.

Liksom \chapref{c2:ch}:s tvåsiffriga display är den här designen byggd av mer än en entitet, med
instansieringen från \secref{c2:sec:submodules} och genericen från \secref{c4:sec:generics}.
Skillnaden är vad delarna gör: displayen instansierade en modul två gånger för att göra samma jobb
på två halvor av sin ingång, medan den här instansierar två \emph{olika} moduler med ett jobb var.
Båda är samma mekanism. Att komponera en design av moduler är inte en särskild teknik för upprepad
logik; det är så designer byggs när de vuxit ur en enda fil.

De två kedjorna bor i två subkomponenter:
\begin{itemize}
  \item \code{reset_sync} är \secref{c4:sec:resetsync}:s tvåvippsresetsynkroniserare. Den tar den
    råa \code{reset_n} och returnerar \code{reset_s2_n}.
  \item \code{button_sync} är \secref{c4:sec:buttonsync}:s trevippssynkroniserare med
    flankdetektor. Den tar den redan synkroniserade \code{reset_s2_n} och pulsar
    \code{button_edge_s2} en gång per tryck.
\end{itemize}

Att hålla dem åtskilda betyder mer än det ser ut. En design utan egen knapp behöver ändå sin reset
synkroniserad, och med två moduler instansierar den helt enkelt den den vill ha, som
\chapref{c8:ch}:s \code{seq_detect_mealy} gör. Hopsvetsade till en modul skulle den designen behöva
instansiera alltihop och koppla bort den halva den inte använder.

\code{button_sync} tar en \emph{vektor} av knappar så att en fil kan betjäna en design med två
knappar. Den här designen har en, så toppnivån deklarerar ett par enelementsvektorer att koppla
till den; det är rördragning, inte logik, två ledningar omdöpta. Fullständig källkod:

\vhdlfile{lectures/L04/led_toggle_sync/led_toggle_sync.vhd}

Toppnivåarkitekturen innehåller ingenting annat än själva växlingsprocessen, och den rör aldrig
annat än de redan synkroniserade \code{reset_s2_n} och \code{button_edge_s2}, aldrig de råa
ingångarna.

De två subkomponenterna finns inte i kursrepot, eftersom att skriva dem är
Övning~\ref{c4:ex:resetsync} och \ref{c4:ex:buttonsync}: \code{reset_sync} är
\secref{c4:sec:resetsync} ovan, och \code{button_sync} är \secref{c4:sec:buttonsync} och
\secref{c4:sec:generics}, båda utskrivna i sin helhet. Kopiera in dina egna i
\file{led_toggle_sync/} för att bygga den.

\paragraph{För hand, i CircuitVerse först}
Innan du rör VHDL, bygg den som ett grindnät: två D-vippor i serie för resetsynkroniseraren; två
för \code{button_n} plus en tredje som lagrar det föregående synkroniserade värdet; en AND-grind
som kombinerar det inverterade nuvarande värdet med det föregående för att producera
\code{button_edge_s2}; och en växlingsvippa för \code{led_s}, med sin inverterade utgång
återkopplad till sin ingång, aktiverad endast när \code{button_edge_s2 = 1}. Sätt klockperioden
till \code{1000 ms} så att du hinner följa varje steg med ögat.

\lecfig[0.95]{lectures/L04/appendix/images/led_toggle_sync_circuit.png}{Synkroniserad
lysdiodsväxlingskrets: dubbelvippsynkroniserare för reset och knapp, flankdetektering, och en
växlingsvippa för lysdioden.}{c4:fig:circuit}

En sak i ritningen behöver läsas noga: vipporna har en \code{preset}-ingång, och på
knappsynkroniserarens tre steg är den bunden till \code{1}. Det är ingen dekoration. I CircuitVerse
är \code{preset} det värde en vippa laddar när dess reset aktiveras, så att binda den till \code{1}
är det som gör att \code{button_s1_n}, \code{button_s2_n} och \code{button_s3_n} kommer ur reset
med värdet \code{1} snarare än \code{0}.

Det stämmer med \code{button_sync}, som resettar samma tre steg till \code{(others => '1')}, och
det spelar roll av samma skäl: knapparna är aktivt låga, så \code{1} betyder ”släppt”. Resetta
kedjan till \code{0} i stället och den startar med påståendet att knappen redan hålls nedtryckt.
Resetvipporna lämnar \code{preset} i fred, eftersom en reset till \code{0} där verkligen är allt
kretsen ber om.

Tryck på knappen i simuleringen och följ hur det fortplantar sig. Lysdioden växlar inte omedelbart:
den väntar på att den synkroniserade, flankdetekterade pulsen ska nå växlingsvippan, vilket tar ett
litet, fast antal klockpulser, precis den latens \secref{c4:sec:sync} kallade priset för säkerheten.

\paragraph{På FPGA:n}
Samma design demonstreras på Terasic DE0-CV (enheten \code{5CEBA4F23C7N}) under passet, med
\code{clock} på \code{50 MHz}-oscillatorn, \code{reset_n} och \code{button_n} på två tryckknappar,
och \code{led} på en lysdiod. Lysdioden växlar vid varje tryck, och enligt
\secref{c4:sec:buttonsync} kan ett riktigt tryck ibland ge mer än en växling på grund av
kontaktstuds.
